\begin{oving}
\begin{itemize}[leftmargin=5.6em,labelwidth=5em,labelsep=0.6em,align=left]
\item[\middels] \href{https://github.com/gustavbv/Matte4C---eksamen/blob/main/oppgavesett/F09-pde-og-bolgelikning/2-Bergmann-eksamener/010-TMA4125-30-2025-kont-oppg6-oppgave.pdf}{Matte 4N kont.\ 2025, oppg.~6} {\color{graa}(\href{https://github.com/gustavbv/Matte4C---eksamen/blob/main/oppgavesett/F09-pde-og-bolgelikning/2-Bergmann-eksamener/010-TMA4125-30-2025-kont-oppg6-svar.pdf}{LF})}: varmelikningen med separasjon
\item[\middels] \href{https://github.com/gustavbv/Matte4C---eksamen/blob/main/oppgavesett/F09-pde-og-bolgelikning/3-hele-oppgaven-i-pensum/022-TMA4130-35-2025-kont-oppg5-oppgave.pdf}{TMA4130 kont.\ 2025, oppg.~5} {\color{graa}(\href{https://github.com/gustavbv/Matte4C---eksamen/blob/main/oppgavesett/F09-pde-og-bolgelikning/3-hele-oppgaven-i-pensum/022-TMA4130-35-2025-kont-oppg5-svar.pdf}{LF})}: isolerte ender (Neumann)
\item[\vanskelig] \href{https://github.com/gustavbv/Matte4C---eksamen/blob/main/oppgavesett/F09-pde-og-bolgelikning/2-Bergmann-eksamener/016-TMA4125-2022-var-oppg8-oppgave.pdf}{Matte 4N vår 2022, oppg.~8} {\color{graa}(\href{https://github.com/gustavbv/Matte4C---eksamen/blob/main/oppgavesett/F09-pde-og-bolgelikning/2-Bergmann-eksamener/016-TMA4125-2022-var-oppg8-svar.pdf}{LF})}: inhomogen varmelikning
\item[\vanskelig] \href{https://github.com/gustavbv/Matte4C---eksamen/blob/main/oppgavesett/F09-pde-og-bolgelikning/3-hele-oppgaven-i-pensum/027-TMA4125-30-35-2024-kont-oppg6-oppgave.pdf}{Matte 4N kont.\ 2024, oppg.~6} {\color{graa}(\href{https://github.com/gustavbv/Matte4C---eksamen/blob/main/oppgavesett/F09-pde-og-bolgelikning/3-hele-oppgaven-i-pensum/027-TMA4125-30-35-2024-kont-oppg6-svar.pdf}{LF})}: Laplace-likningen på et kvadrat med separasjon
\end{itemize}
\end{oving}
